\documentclass{article}
\usepackage[UTF8]{ctex}
\usepackage{amsmath,amssymb,amsfonts}
\usepackage{enumitem}
\usepackage{geometry}
\geometry{a4paper, margin=1in}
% === 错题本排版优化（不参与题目标识） ===
\setlength{\parindent}{0pt}
\setlength{\parskip}{0.35em plus 0.10em minus 0.08em}
\setlength{\emergencystretch}{3em}
\setlist{itemsep=0.25em, topsep=0.40em, parsep=0pt}
\allowdisplaybreaks[2]
% === 错题本排版优化结束 ===
\begin{document}
\section*{云南大学2025年高等代数真题}
\section*{184．云南大学2025年高等代数试卷}
一．填空题．\\

1．设 \(\alpha, \beta\) 是 \(n\) 维欧氏空间 \(V\) 中的两个向量，且 \(|\alpha|=2,|\beta|=3\) ．则 \(|\alpha+\beta|^{2}+|\alpha-\beta|^{2}=\) \(\_\_\_\_\)．

2．设 \(V_{1}, V_{2}\) 是有限维线性空间 \(V\) 的两个子空间，若维 \(\left(V_{1}\right) \leq\) 维 \(\left(V_{2}\right)\) ，维 \(\left(V_{1}+V_{2}\right)=2025\) ，维\(\left(V_{1} \cap V_{2}\right)=2024\) ，则维 \(\left(V_{1}\right)=\) \(\_\_\_\_\)．

3．实二次型 \(-4 x_{1} x_{2}+2 x_{1} x_{3}+2 x_{2} x_{3}\) 的规范型为 \(\_\_\_\_\)．

4．若 \(A\) 是 \(n\) 阶可对角化幂零阵，则 \(A=\) \(\_\_\_\_\)．

5．若2024 阶方阵 \(A=\left(a_{i j}\right)\) ，其中 \(a_{i j}=\left\{\begin{array}{ll}0, & \text { 若 } j-1>i ; \\ C_{i}^{j-1}, & \text { 若 } j-1 \leq i .\end{array}\right.\) 这里 \(C_{i}^{j-1}=\frac{i!}{(j-1)!(i-j+1)!}\) ，则 \(|A|=\)\\
\(\_\_\_\_\)．

6．设 \(f(x)=x^{3}-x+2, g(x)=x^{4}+x-1\) ，则 \((f(x), g(x))=\) \(\_\_\_\_\)．

二．求多项式 \(f(x)=x^{5}-5 x^{4}+7 x^{3}-2 x^{2}+4 x-8\) 的所有有理根，计算重数．

三。设 10 阶方阵 \(A=\left(a_{i j}\right)\), 其中 \(a_{i j}=\left\{\begin{array}{ll}1, & \text { 若 } i=j+1 ; \\ 0, & \text { 若 } i \neq j+1 .\end{array}\right.\) 求 \(A^{3}\) 的 Jordan 标准型。

四．记 \(M_{2}(\mathbb{R})\) 为由所有2阶实方阵构成的欧氏空间，内积定义为：\((A, B)=\operatorname{tr}\left(A^{\mathrm{T}} B\right)\) ，对 \(\forall A, B \in\) \(M_{2}(\mathbb{R})\) ，这里 \(\operatorname{tr}(H)\) 是指方阵 \(H\) 的所有主对角元的和。

\begin{itemize}

  \item[1．] 记 \(W=\left\{A \in M_{2}(\mathbb{R}) \mid \operatorname{tr}(A)=0\right\}\) ，求 \(W\) 的维数与 \(W^{\perp}\) 。

  \item[2．] 设 \(A=\left(\begin{array}{ll}2 & 1 \\ 1 & 2\end{array}\right)\) ，求 \(A\) 在 \(W\) 上的内投影（亦称正交投影）。

  \item[3．] 对 \(A=\left(\begin{array}{ll}2 & 1 \\ 1 & 2\end{array}\right), B=\left(\begin{array}{ll}1 & 1 \\ 1 & 1\end{array}\right), C=\left(\begin{array}{ll}0 & 1 \\ 1 & 0\end{array}\right)\) ，定义 \(M_{2}(\mathbb{R})\) 上的线性变换 \(\sigma\) 为
\[
\sigma(D)=A D+B D C, \forall D \in M_{2}(\mathbb{R}) .
\]
\end{itemize}

请找出 \(M_{2}(\mathbb{R})\) 的一组标准正交基，使得 \(\sigma\) 在此基下的矩阵是对角阵，并写出此对角阵。

五．证明：反对称矩阵的秩一定是偶数．

六．设 \(V\) 是数域 \(F\) 上的一个 \(n\) 维线性空间，\(\sigma\) 是 \(V\) 上的线性变换．对多项式 \(f(x), g(x) \in F[x]\) ，记 \(d(x)=(f(x), g(x)), m(x)=[f(x), g(x)]\) ，对多项式 \(h(x) \in F[x]\) ，记 \(h(\sigma)(V)=\{h(\sigma) \alpha \mid \alpha \in V\) ．\} 为线性变换 \(h(\sigma)\) 的象空间．证明：

\begin{itemize}

  \item[1．] \(m(\sigma)(V)=f(\sigma)(V) \cap g(\sigma)(V)\) ．

  \item[2．] \(d(\sigma)(V)=f(\sigma)(V)+g(\sigma)(V)\) ．

  \item[3．] 存在多项式 \(h_{1}(x), h_{2}(x) \in F[x]\) ，使得 \(h_{1}(\sigma)(V)\) 与 \(h_{2}(\sigma)(V)\) 都是 \(\sigma\) 的不变子空间，且 \(\sigma \mid\) \(h_{1}(\sigma)(V)\) 可逆，\(\sigma \mid h_{2}(\sigma)(V)\) 幂零，且 \(V=h_{1}(\sigma)(V) \oplus h_{2}(\sigma)(V)\) ．
\end{itemize}

七．设 \(\varepsilon_{i}=(0, \cdots, 0,1,0, \cdots, 0) \in \mathbb{R}^{n+1}\) 第 \(i\) 个元素为 1 ，其余皆为 \(0, i=1,2, \cdots, n+1\) ，这里

\[
\mathbb{R}^{n+1}=\left\{\left(a_{1}, a_{2}, \cdots, a_{n+1}\right) \mid a_{i} \in \mathbb{R}, i=1,2, \cdots, n+1\right\}
\]

是 \(n+1\) 维欧氏空间，其上内积为通常内积。对 \(i=1,2, \cdots, n\) ，记 \(\alpha_{i}=\varepsilon_{i}-\varepsilon_{i+1}\) ，又记 \(\mathbb{R}^{n+1}\) 上的 \(n\)个线性变换 \(\sigma_{i}\) 的定义为 \(\sigma_{i}(\beta)=\beta-\left(\beta, \alpha_{i}\right) \alpha_{i}, \forall \beta \in \mathbb{R}^{n+1}\) ，设

\[
W=\left\{\sigma_{i_{1}} \sigma_{i_{2}} \cdots \sigma_{i_{t}} \mid t \text { 取遍自然数, 对任意的 } i_{1}, i_{2}, \cdots, i_{t} \in\{1,2, \cdots, n\}\right\} \text {. }
\]

证明：

\begin{itemize}

  \item[1．] 线性变换 \(\sigma_{i}\) 是第二类正交变换。

  \item[2．] 对任意的 \(\gamma \in \mathbb{R}^{n+1}\) ，存在 \(\sigma \in W\) ，使得 \(\left(\sigma(\gamma), \alpha_{i}\right) \geq 0\) 对一切 \(i=1,2, \cdots, n\) 成立，并且满足该不等式的向量 \(\sigma(\gamma)\) 由 \(\gamma\) 唯一确定．
\end{itemize}
\end{document}
